% French back cover: page preceding the English back cover, with title,
% author, French abstract and keywords; no page number, running head or line
% numbers. Always on a recto when printed recto-verso, and its margins are
% centered like those of the title page.
\cleartorightpage
\centeredmargins
% PDF bookmark only (not listed in the table of contents).
\phantomsection
\pdfbookmark[0]{Back Cover}{backcover}
\thispagestyle{empty}
\nolinenumbers
\label{chap:french_back_cover}
\begin{center}
  {\LARGE\bfseries \ThesisTitle\par}
  \medskip
  {\large \ThesisAuthor\par}
\end{center}
\bigskip
\noindent\textbf{Résumé}

Les agents d'apprentissage par renforcement profond atteignent de hautes performances sur un large éventail de tâches.
Cependant, les réseaux de neurones qui implémentent leurs politiques sont difficiles à interpréter en raison de leur nature connexionniste.
Cette difficulté vient de ce que l'information est distribuée dans l'ensemble du réseau, ce qui rend difficile de relier ses composants internes aux décisions de l'agent.

Dans cette thèse, l'explicabilité de ces agents est obtenue par l'analyse de leurs espaces latents.
Cette démarche s'inscrit dans une perspective mécaniste, qui cherche à expliquer les décisions des agents en analysant les représentations internes apprises par leurs réseaux de neurones.

La thèse commence par passer en revue les méthodes de l'état de l'art pour l'analyse des espaces latents, afin de donner une vue d'ensemble du domaine.
Elle précise ensuite ce qui constitue une explication dans le contexte de l'apprentissage par renforcement explicable.
Sur ces bases, elle passe en revue l'état de l'art de l'apprentissage par renforcement explicable à l'aide d'une taxonomie spécifiquement adaptée à l'apprentissage par renforcement.

En s'appuyant sur ces deux revues, la thèse introduit \textit{Clustering Agent Representation}, une méthode d'explicabilité qui analyse l'espace latent des politiques d'apprentissage par renforcement afin d'identifier des structures significatives dans leurs représentations internes.

Pour évaluer objectivement les structures obtenues, trois métriques sont proposées.
\textit{Clustering Universality} mesure leur cohérence entre des politiques entraînées indépendamment, ce qui établit que les structures identifiées ne sont pas dues au hasard.
\textit{Abstraction Score} mesure leur écart à la structure des observations, à celle des actions et à celle d'un réseau non entraîné, ce qui montre que ces structures émergent de l'apprentissage.
\textit{Adjusted Clustering Universality} combine ces deux mesures pour fournir un critère de choix des représentations à analyser.

Les expériences menées sur six environnements Atari montrent que des politiques entraînées indépendamment convergent vers des structures de clustering similaires.
Comme cette similarité ne s'explique ni par les observations, ni par les actions, ni par l'architecture du réseau, elle indique une convergence des concepts appris par ces politiques.
Enfin, une analyse qualitative montre que les clusters obtenus correspondent à des aspects interprétables de la tâche.

Ces résultats montrent que l'analyse des espaces latents est une voie prometteuse pour expliquer les agents d'apprentissage par renforcement, et ouvrent de nouvelles perspectives pour leur compréhension mécaniste.

\bigskip
\noindent\textbf{Mots-clés :} Intelligence Artificielle Explicable, Apprentissage Profond, Apprentissage par Renforcement, Explicabilité, Explicabilité Mécanistique, Interprétabilité Mécanistique, Espace Latent, Clustering, Regroupement.
